\section{Even classes and triangles of auxiliary places}\label{sec:triangles}

Ribet's lemma (\cref{thm:ribet-lemma}) needs, for each character $\mathrm{cyc}^{-e}$ of the string, a triangle cost for the modules of character cocycles. These are cocycles in the
Tate twists $(O/\ell^\nu)(-e)$; their \emph{even} parts are controlled by auxiliary places in triangular position, chosen with Kummer maps over the real subfield, and their
\emph{odd} parts by a class number. No bound for the even classes themselves is needed: the reciprocity laws of the Kummer maps remove them.

\subsection{Even and odd classes}\label{ss:even-classes}
Let $E\subset\Omega$ be a CM field containing a primitive $5040$th root of unity, $F$ its maximal real subfield, $O$ a coefficient ring (the integers of a finite Galois extension of $\Q_\ell$ in $\bar\Q_\ell$), and
$G_\nu\subseteq G$ the elements fixing the $\ell^\nu$th roots of unity. For a set $S$ of places, the \emph{even classes} of exponent $e$ are the cocycles $c\colon G\to(O/\ell^\nu)(-e)$, unramified outside $\ell$, $S$ and $\bar S$, with
$c(g^c)=c(g)$, vanishing on the elements of trivial cyclotomic character in the decomposition groups above $\ell$; the \emph{odd classes} are the cocycles with the same local conditions, unramified off $\ell$, that
change sign under $g\mapsto g^c$. A family of places $w_j$ off $\ell$ with the chosen tame generators $\tau_j$ and elements $\varphi_j\in G_\nu$ is a \emph{tame generator family}. A \emph{triangle of primes} for an element $a$ is such a family,
indexed by a finite linear order $J$ with monotone exponents $0<\mathrm{ty}(j)<t$, such that: a finite set $s$ of test points of trivial cyclotomic character, with $(\#J+\#s+1)h\le c$, detects together with $a$ the
cocycles unramified off $\ell$, and for each exponent and each $g\in s$ the even classes unramified off $\ell$ take at some $\varphi_j$ of that exponent twice their value at $g$; an even class of exponent $\mathrm{ty}(j)$ that kills the
earlier $\tau_i$ and the earlier $\varphi_i$ of the same exponent has its value at $\tau_j$ killed by $\ell^h$; and $a$ separates the powers $\mathrm{cyc}(a)^{-e}$, $e<t$, up to $\ell^h$.
Moreover $4(\mathrm{cyc}(a)^{-e}-1)^2$ divides $\ell^h$ for $0<e<t$. The \emph{local set} above $\ell$ is the union of the inertia groups above $\ell$ if $e=0$, and otherwise the set of the elements of trivial cyclotomic character in the
decomposition groups above $\ell$; the even and odd classes vanish on it, and so does every cocycle that obeys the crystalline relations above $\ell$ (\cref{ss:B-padic}).
\lean{New.Selmer.evenClasses, New.Selmer.IsTameGeneratorFamily, New.Selmer.withConjInv, New.Selmer.withConjInv_true, New.Selmer.withConjInv_false, New.Selmer.IsPrimeTriangle, New.Selmer.IsAuxTwist, New.Selmer.four_dvd_cyclotomicChar_sub_one, New.Selmer.tateTwist_eq_one_of_cyclotomicChar_eq_one, New.Selmer.IsTameGeneratorFamily.cyclotomicInv_withConjInv_eq_one_and_dvd, New.Selmer.IsTameGeneratorFamily.mem_selmerModule_of_tame_eq_zero, New.Selmer.eq_zero_of_forall_mem_Z1_of_mem_evenClasses}

\begin{lemma}[Even parts and tame generators]\label{lem:even-parts}
\begin{enumerate}[label=\textup{(\alph*)}]
\item For $e\ne0$ the even part $c+c\circ(g\mapsto g^c)$ of a cocycle with values in $(O/\ell^\nu)(-e)$, unramified outside $\ell$, $S$, $\bar S$ and zero on the local set above $\ell$, is an even class.
\item A cocycle unramified off $S$ and the $w_j$, $\bar w_j$ that vanishes at the $\tau_j$ and at the $(\tau_j^c)^{-1}$ is unramified off $S$.
\item An element $b$ of $O/\ell^\nu$ killed by every additive map $u$ to $\Z/\ell^\nu$ for which $u\circ c$ is a cocycle, for an even class $c$, is zero.
\end{enumerate}
\end{lemma}

\begin{proof}
(a) Conjugation $g\mapsto g^c$ is an automorphism of $G$ that fixes the cyclotomic character, $\mathrm{cyc}(cgc)=\mathrm{cyc}(g)$, so $c\circ(g\mapsto g^c)$ is again a cocycle for $\mathrm{cyc}^{-e}$: $c((gh)^c)=c(g^c)+\mathrm{cyc}(g)^{-e}c(h^c)$. It carries the inertia and decomposition groups
above a place to those above the conjugate place, which is above $\ell$, in $S\cup\bar S$, or neither, as the original one; so it keeps the ramification conditions and, the cyclotomic character being fixed, the vanishing on the
local set (for $e\ne0$ this set does not depend on $e$). The sum $c+c\circ(g\mapsto g^c)$ is invariant under the involution $g\mapsto g^c$.

(b) At a place $A$ above $w_j$, which is not above $\ell$, the cyclotomic character, hence the twist, is trivial on inertia, so the cocycle restricted to the inertia group is a homomorphism to an $\ell$-power torsion group. It factors
through the tame quotient, wild inertia being pro-$p$ with $p\ne\ell$, and the $\ell$-part of the tame quotient is procyclic and topologically generated by the image of the chosen generator $\tau_j$. Hence it vanishes on the inertia
group at $A$. For another place $sA$ over $w_j$ and $g$ in its inertia group, $z(g)=z(s\cdot s^{-1}gs\cdot s^{-1})$ is $s\cdot z(s^{-1}gs)$ up to terms that cancel because $g$ acts as the identity, so $z$ vanishes there too. At $\bar w_j$ the same holds with the tame
generator $(\tau_j^c)^{-1}$ at the conjugate place $cA$.

(c) Characters with values in $\Q/\Z$ separate the points of $O/\ell^\nu$, and on an $\ell^\nu$-torsion group each of them factors through $\Z/\ell^\nu$; so $b=0$ if every additive $u\colon O/\ell^\nu\to\Z/\ell^\nu$ kills it. For each such $u$, $u\circ c$ is a continuous
cocycle for $\mathrm{cyc}^{-e}$ modulo $\ell^\nu$, as the twist acts by integers modulo $\ell^\nu$ and $u$ is additive.
\end{proof}

\subsection{Kummer maps}\label{ss:kummer-maps}
A Frobenius datum $x$ (a prime of degree one, a place of $\Omega$ above it and a Frobenius element) is \emph{$\nu$-split} if its prime is not above $\ell$ and its Frobenius element lies in $G_\nu$. For a finite family $x$ of $\nu$-split places,
a \emph{Kummer map} of exponent $e$ off a finite set $T$ of primes of $F$ is the restriction to $G$ of a relation over $F$ for the module $\Z/\ell^\nu$ with the action of $(\mathrm{cyc}^+)^{-e}$, $\mathrm{cyc}^+$ the extension of the cyclotomic
character to $\Gal(\Omega/F)$ that is $1$ at $c$: a locally constant equivariant homomorphism $y$ from $G_\nu$ to the Cartier dual, zero on inertia off $T$, satisfying reciprocity $\sum_py(\varphi_p)(z(\gamma_p))=0$ against the
cocycles $z$ over $F$ with the local conditions above $\ell$ and below $x$, for every finite family of distinct primes $p\notin T$.
\lean{New.Selmer.IsNuSplitPlace, New.Selmer.primeBelowPlus, New.Selmer.primeBelowPlus_star, New.Selmer.q_primeBelowPlus, New.Selmer.toGalPlus, New.Selmer.conjPlus, New.Selmer.isOpenEmbedding_toGalPlus, New.Selmer.index_range_toGalPlus, New.Selmer.isComplexConj_conjPlus, New.Selmer.conjPlus_mul_self, New.Selmer.isFrob_toGalPlus_iff, New.Selmer.toGalPlus_mem_inertiaGroup_iff, New.Selmer.cyclotomicChar_toGalPlus, New.Selmer.toGalPlus_conjGal, New.Selmer.notMem_range_toGalPlus, New.Selmer.exists_toGalPlus_eq_ite, New.Selmer.ofGalPlus, New.Selmer.toGalPlus_ofGalPlus, New.Selmer.toGalPlus_ofGalPlus_of_mem, New.Selmer.toGalPlus_ofGalPlus_of_notMem, New.Selmer.continuous_ofGalPlus, New.Selmer.ofGalPlus_mul_eq_or, New.Selmer.inertiaGroup_le_range_toGalPlus, New.Selmer.ofGalPlus_mem_inertiaAbove, New.Selmer.cyclotomicPlus, New.Selmer.cyclotomicPlus_apply, New.Selmer.continuous_cyclotomicPlus, New.Selmer.cyclotomicPlus_eq_one_of_mem_inertiaAbove, New.Selmer.cyclotomicPlus_eq_one_of_isComplexConj, New.Selmer.cyclotomicKerPlus, New.Selmer.cyclotomicPlusMod, New.Selmer.cyclotomicPlusMod_apply, New.Selmer.cyclotomicKerPlus_le_ker, New.Selmer.isOpen_ker_cyclotomicPlusMod, New.Selmer.tateTwistPlus, New.Selmer.tateTwistPlus_ρ, New.Selmer.tateTwistPlus_ρ_eq_one, New.Selmer.cyclotomicPlusMod_eq_one, New.Selmer.familyPlace, New.Selmer.exists_familyPlace_eq, New.Selmer.isPlaceOver_familyPlace, New.Selmer.exists_isFrob_familyPlace, New.Selmer.familyPrimes, New.Selmer.mem_familyPrimes, New.Selmer.cyclotomicPlusMod_eq_one_of_mem_decompGroup, New.Selmer.kummerSetting, New.Selmer.IsKummerMap, New.Selmer.HasKummerType, New.Selmer.HasKummerType.map_one, New.Selmer.HasKummerType.map_inv, New.Selmer.HasKummerType.map_pow, New.Selmer.IsKummerMap.hasKummerType, New.Selmer.comp_ofGalPlus_mem_Z1, New.Selmer.two_nsmul_comp_ofGalPlus_mem_cocycles, New.Selmer.primeBelowPlus_mem_of_mem_conjClosure, New.Selmer.IsKummerMap.sum_eq_zero, New.Selmer.exists_isTameGeneratorFor_toGalPlus, New.Selmer.exists_apply_eq_nsmul_of_mem_evenClasses, New.Selmer.IsKummerMap.nsmul_apply_tameGenerator_eq_zero, New.Selmer.reductionModL, New.Selmer.kummerReductions, New.Selmer.natCast_dvd_sub_of_cast_eq, New.Selmer.exists_apply_toGalPlus_eq_ofZMod, New.Selmer.finite_kummerReductions, New.Selmer.card_reductions_le, New.Selmer.exists_card_lt_finrank_kummerReductions}

Write $G_F=\Gal(\Omega/F)\supseteq G$ and $\mathrm{cyc}^+$ for the character of $G_F$ that is $\mathrm{cyc}$ on $G$ and $1$ at $c$: the cyclotomic character times the quadratic character of $E/F$. The map $G_F\to G$ sending $g$ and $gc$ to $g$,
$g\in G$, is continuous, and multiplicative up to replacing the second factor $s$ by $s^c$. The extension $E/F$ is unramified at the finite places: otherwise an element of inertia would act on $E$ as complex conjugation, but there is an
integer $y$ of $E$ with $y-\bar y$ a unit (from the roots of unity of order $5040$ in $E$), which would lie in the prime below the place. A \emph{tame generator normalised over $F$} at a prime of degree one of $E$ is an element of inertia in $G$
that moves an $m$-th root of a uniformiser of $F$ by the fixed root of unity; such generators exist.

\begin{proposition}[Kummer maps]\label{prop:kummer-maps}
\begin{enumerate}[label=\textup{(\alph*)}]
\item A Kummer map $\psi$ is locally constant; on $G_\nu$ it is a homomorphism negated by complex conjugation, and conjugation by $g$ multiplies it by $\mathrm{cyc}(g)^{e+1}$.
\item (Reciprocity.) For $\nu$-split places $w_j$ above distinct primes of $F$ outside $T$ and an even class $z$ of exponent $e$ for $w$ vanishing at the Frobenius elements of $x$, $\sum_j\psi(\varphi_j)\cdot2z(\gamma_j)=0$. In particular,
if $\psi$ vanishes at the Frobenius elements of $w$ except at $w_j$, every multiple of $2\psi(\varphi_j)$ kills $z(\tau_j)$.
\item If $E$ has at least four infinite places, for every finite family of $\nu$-split places above distinct primes of $F$ there is $T$ such that the reductions modulo $\ell$ of the Kummer maps span a space of dimension
larger than the number of places.
\end{enumerate}
\end{proposition}

\begin{proof}
(a) The relation $y$ of $\psi$ is a locally constant homomorphism on $G_\nu$, and $y(g)(a)=a\psi(g)/\ell^\nu$ determines $\psi$, both sides being additive in $a$ and agreeing at $1$. By the transformation law of a relation,
$y(gug^{-1})=\mathrm{cyc}(g)\cdot y(u)\circ g^{-1}$, and $g^{-1}$ acts on $\Z/\ell^\nu$ by $\mathrm{cyc}^+(g)^e$; so conjugation by $g\in G_F$ multiplies $\psi$ by $\mathrm{cyc}(g)\mathrm{cyc}^+(g)^e$. For $g\in G$ this is $\mathrm{cyc}(g)^{e+1}$, and for $g=c$ it is $\mathrm{cyc}(c)=-1$.

(b) \emph{Extension to $F$.} Replacing $z$ by $u\circ z$ for the additive maps $u$ to $\Z/\ell^\nu$ (\cref{lem:even-parts}(c)) we may assume $z$ has values in $\Z/\ell^\nu$. Put $z^+(g)=z^+(gc)=z(g)$ for $g\in G$. As $z$ is invariant under $g\mapsto g^c$, $z^+$ is a cocycle of
$G_F$ for $(\mathrm{cyc}^+)^{-e}$, and $2z^+$ lies in the cocycles of the Kummer setting for the primes $P$ below the $w_j$: it vanishes at complex conjugations; it is unramified off $S_0\cup P$, as $E/F$ is unramified at the finite places;
above $\ell$, for $g$ in the decomposition group with $\mathrm{cyc}^+(g)=1$, $2z^+(g)=z^+(g^2)$ with $g^2$ in $G$, in the decomposition group and in $\ker\mathrm{cyc}$, where $z$ vanishes; and below $x$, the decomposition group is topologically
generated by inertia and a Frobenius element and lies in $\ker\mathrm{cyc}^+$ modulo $\ell^\nu$, $x$ being $\nu$-split, so $2z^+$ restricted to it is a homomorphism vanishing at the Frobenius element (by hypothesis) and on inertia ($z$ being
unramified there). \emph{Reciprocity.} The primes of $F$ below the $w_j$ split in $E$, with the arithmetic Frobenius elements $\varphi_j^{-1}$, and the $\gamma_j$ are tame generators normalised over $F$. The reciprocity of the relation gives
$\sum_jy(\varphi_j^{-1})(2z^+(\gamma_j))=0$, that is $-\sum_j\psi(\varphi_j)\cdot2z(\gamma_j)/\ell^\nu=0$. \emph{One term.} If $\psi(\varphi_i)=0$ for $i\ne j$, then $2\psi(\varphi_j)z(\gamma_j)=0$. On the inertia group at $w_j$ the cocycle $z$ is a homomorphism
through the procyclic $\ell$-part of tame inertia (\cref{lem:even-parts}(b)), so $z(\tau_j)$ is a multiple of $z(\gamma_j)$; hence $n=2\psi(\varphi_j)a$ kills $z(\tau_j)$.

(c) Apply \cref{thm:B-relations-mod-l} over $F$ to $M=\Z/\ell^\nu$ with the action $(\mathrm{cyc}^+)^{-e}$, $U=G_\nu$, $\kappa=\mathrm{cyc}^+$ and the local conditions of the Kummer setting. The hypotheses hold: $M[\ell]$ has $\ell$ elements; $\mathrm{cyc}^+$ takes
values in $1+4\Z_\ell$, since $E$ contains the fourth roots of unity and $\mathrm{cyc}^+(c)=1$; its kernel lies in $G_\nu\cup G_\nu c$; and $F$ has as many infinite places as $E$, at least four. Hence for some $T\supseteq S_0$ the reductions of the relations
off $T$ number at least $\ell^{k+1}$, with $k$ at least the number of the given places, their primes below being distinct and off $\ell$. A relation is determined by its Kummer map, and its reduction by the reduction of the Kummer map,
so the reductions of the Kummer maps span a space with at least $\ell^{k+1}$ elements, of dimension at least $k+1$.
\end{proof}

\subsection{Triangles of primes}\label{ss:prime-triangles}
\lean{New.Selmer.mul_prod_pow_sub_pow_dvd, New.Selmer.linearIndependent_of_notMem_span_Iio, New.Selmer.surjective_of_linearIndependent_cast, New.Selmer.exists_isKummerMap_notMem_span, New.Selmer.exists_apply_eq_prod_mul, New.Selmer.exists_separating_element, New.Selmer.additive_apply_pow_mul_conj, New.Selmer.HasKummerType.eventually_eq_two_mul, New.Selmer.evenClasses_eventually_eq_two_nsmul, New.Selmer.exists_forall_of_wellFoundedLT, New.Selmer.TriangleEntry, New.Selmer.TriangleEntry.MayFollow, New.Selmer.eventually_primeBelowPlus_notMem, New.Selmer.TriangleEntry.exists_mayFollow, New.Selmer.TriangleEntry.nonempty, New.Selmer.TriangleEntry.exists_forall_mayFollow, New.Selmer.TriangleEntry.primeBelowPlus_injective_of_mayFollow, New.Selmer.TriangleEntry.isTameGeneratorFamily, New.Selmer.TriangleEntry.nsmul_tameGenerator_eq_zero_of_mayFollow, New.Selmer.IsNuSplitPlace.prime_pow_dvd_q_sub_one, New.Selmer.isConjApart_of_injective_primeBelowPlus, New.Selmer.exists_isPrimeTriangle}
\lean{New.Selmer.add_conjGal_mem_evenClasses, New.Selmer.apply_withConjInv_false_eq_neg, New.Selmer.apply_eq_zero_of_mem_obeying_crysRelations, New.Selmer.relLength_torsionBy_le_length_shaOdd, New.Selmer.hasTriangleCost_tateTwist}

\begin{lemma}[Choosing the next place]\label{lem:next-place}
Let $\nu>0$, $t+2\le y$, $\mathrm{cyc}(\gamma)^a-1\mid\ell^h$ for $0<a<t$, $4\mid\ell^{3h}$, and let $W$ contain most Frobenii near $\delta^{\ell^\nu}g'$ for some $\delta$ of trivial cyclotomic character, for every $g'$. Let finitely many places above distinct primes
of $F$ carry maps $\psi_i$ of Kummer type of exponents $\mathrm{ty}(i)$, those of exponent $e$ independent modulo $\ell$, and let $g$ have trivial cyclotomic character.
\begin{enumerate}[label=\textup{(\alph*)}]
\item There is $u\in G_\nu$ with $\psi_i(u)=0$ for all $i$ and $\psi_0(u)=\prod_{k\ne e}(\mathrm{cyc}(\gamma)^e-\mathrm{cyc}(\gamma)^k)$, for a new Kummer map $\psi_0$ of exponent $e$ for the earlier places of exponent $e$ whose reduction is outside the span of theirs.
\item There is a new place $x_0$ with Frobenius datum in $W$, above a new prime of $F$ avoiding the primes excluded by all the maps, such that the earlier maps vanish at its Frobenius element $\varphi$, $2\psi_0(\varphi)$ divides $\ell^{yh}$, and
every even class of exponent $e$ unramified off $\ell$ takes at $\varphi$ twice its value at $g$.
\item In a family of places chosen in this way along a finite linear order, an even class of exponent $\mathrm{ty}(j)$ that vanishes at the earlier $\tau_i$ and at the earlier $\varphi_i$ of the same exponent has its value at $\tau_j$ killed by
$\ell^{yh}$.
\end{enumerate}
\end{lemma}

\begin{proof}
(a) By \cref{prop:kummer-maps}(c) the reductions of the Kummer maps of exponent $e$ for the earlier places of exponent $e$, off a suitable $T_0$, span more dimensions than there are earlier maps, so some $\psi_0$ has its reduction outside
their span. The maps of exponent $e$ together with $\psi_0$ are homomorphisms on $G_\nu$ with independent reductions, so by Nakayama's lemma they are jointly onto: some $g_1\in G_\nu$ has $\psi_0(g_1)=1$ and $\psi_i(g_1)=0$ for $\mathrm{ty}(i)=e$. For
$c=\mathrm{cyc}(\gamma)$ and a map $\psi$ of exponent $d$, conjugation by $\gamma$ followed by the power $c^{-1}$ (an integer modulo $\ell^\nu$) multiplies $\psi$ by $c^{-1}c^{d+1}=c^d$; combining such operations as products and quotients in $G_\nu$, some $u$
has $\psi(u)=\prod_{k\ne e}(c^d-c^k)\psi(g_1)$ for every map $\psi$ of exponent $d$. This vanishes for $d\ne e$, and for $d=e$ it is the product times $\psi(g_1)$.

(b) \emph{The target.} Let $\sigma=\delta^{\ell^\nu}\,g\,g^c\,u\,(u^c)^{-1}$, with $\delta$ given by the hypothesis on $W$ for $g'=gg^cu(u^c)^{-1}$; it lies in $G_\nu$. For a map $\psi$ of Kummer type, a homomorphism on $G_\nu$ negated by complex conjugation,
$\psi(\sigma)=\ell^\nu\psi(\delta)+\psi(g)-\psi(g)+\psi(u)+\psi(u)=2\psi(u)$, as $g$, $\delta$ have trivial cyclotomic character and so lie in $G_\nu$. For an even class $z$ of exponent $e$ unramified off $\ell$, a homomorphism on $G_\nu$ with values killed by $\ell^\nu$
and invariant under $g\mapsto g^c$, $z(\sigma)=z(g)+z(g)+z(u)-z(u)=2z(g)$. These identities hold near $\sigma$: the maps are locally constant, and there are finitely many even classes unramified off $\ell$ with values in $\Z/\ell^\nu$, to
which the general ones reduce by \cref{lem:even-parts}(c). \emph{Chebotarev.} The filter of most Frobenii near $\sigma$ is proper (\cref{thm:B-classfield}(d)); so some Frobenius datum $x_0\in W$ lies near $\sigma$ and avoids the finitely many
primes above $\ell$, below the earlier places, and in the sets excluded by the maps. Its Frobenius element $\varphi$ has $\psi_i(\varphi)=2\psi_i(u)=0$, $\psi_0(\varphi)=2\psi_0(u)$ and $z(\varphi)=2z(g)$, and $x_0$ is $\nu$-split, $\varphi$ being near $\sigma\in G_\nu$. \emph{The divisibility.}
$2\psi_0(\varphi)=4\prod_{k\ne e}(c^e-c^k)$, and $c^e-c^k=c^k(c^{e-k}-1)$ or $-c^e(c^{k-e}-1)$ divides $\ell^h$; as $4\mid\ell^{3h}$ and there are $t-1$ factors, $2\psi_0(\varphi)$ divides $\ell^{3h}\ell^{(t-1)h}=\ell^{(t+2)h}$, which divides $\ell^{yh}$ as $t+2\le y$.

(c) Let $z$ be such an even class, of exponent $e=\mathrm{ty}(j)$. It vanishes at the earlier $\tau_i$, hence, being even, at the $(\tau_i^c)^{-1}$, so it is unramified at the earlier places (\cref{lem:even-parts}(b)): an even class for the places from $j$
on. The $j$-th map $\psi_j$ is a Kummer map for the earlier places of exponent $e$, at whose Frobenius elements $z$ vanishes, and it vanishes at the Frobenius elements of the later places, chosen after it; the places from $j$ on lie
above distinct primes of $F$ outside its excluded set. By \cref{prop:kummer-maps}(b), every multiple of $2\psi_j(\varphi_j)$ kills $z(\tau_j)$, and $\ell^{yh}$ is one by (b).
\end{proof}

\begin{theorem}[Triangles exist]\label{thm:triangles}
Let $E$ have at least four infinite places and lie in the cyclotomic field of level $N_E$, let $N_E+t<y$, let $\gamma$ generate $G$ modulo the kernel of $\mathrm{cyc}$, let $\mathrm{cyc}(\gamma)^a-1$ divide $\ell^h$ for $0<a<t$, let $\ell\le y^2$ or
$h=0$, and let a set $W$ of Frobenius data contain most Frobenii near $\delta^{\ell^\nu}g$ for some $\delta$ of trivial cyclotomic character, for all $\nu$ and $g$. Then there is a finite linear order $J$ with a monotone exponent map such that for
each $\nu>0$ some places of level $\ell^\nu$ indexed by $J$ have their Frobenius data in $W$ and form a triangle of primes for $\gamma^{-1}$ with $h$ replaced by $yh$ and $c=y^{22}h$.
\end{theorem}

\begin{proof}
\emph{Test points.} By \cref{thm:B-detecting}(a) for $g_0=\gamma^{-1}$ there is a finite set $s$ of elements of trivial cyclotomic character, with $\#s\le y^{20}$ when $N_E\le y$ and $\ell\le y^2$, such that a cocycle in a Tate twist modulo $\ell^n$ unramified off $\ell$
that vanishes at $\gamma^{-1}$ and on $s$ is zero. \emph{The order.} Let $J$ be the set of the pairs $(e,g)$ with $0<e<t$ and $g\in s$, ordered lexicographically with the exponent first (after well-ordering $s$), and $\mathrm{ty}(e,g)=e$, which is
monotone. \emph{The choice.} Fix $\nu>0$. Since $4\mid\mathrm{cyc}(\gamma)^e-1\mid\ell^h$ for $e>0$ (the cyclotomic character is $\equiv1$ modulo $4$ on $G$), $4(\mathrm{cyc}(\gamma)^e-1)^2$ divides $\ell^{3h}$, and $4\mid\ell^{3h}$ unless $J$ is empty. By well-founded recursion along $J$, choose at
each $j=(e,g)$ a map and a place by \cref{lem:next-place}(a), (b) with the test element $g$; by induction the earlier maps of exponent $e$ are independent modulo $\ell$, a family none of whose members lies in the span of the earlier
ones being independent. \emph{The triangle.} The places lie above distinct primes of $F$, so they are distinct from each other and from their conjugates, and are $\nu$-split, so their norms are $\equiv1$ modulo $\ell^\nu$: they are
Taylor--Wiles places of level $\ell^\nu$, with Frobenius data in $W$. The conditions of a triangle of primes for $\gamma^{-1}$, with $yh$ in place of $h$ and $c=y^{22}h$: the detection is that of the test points; at the place $(e,g)$ the even
classes unramified off $\ell$ take twice their value at $g$ (\cref{lem:next-place}(b)); the triangle condition is \cref{lem:next-place}(c); $\prod_{k\ne e}(u^e-u^k)$, $u=\mathrm{cyc}(\gamma)$, divides $\ell^{(t-1)h}\mid\ell^{yh}$, and $4(u^e-1)^2\mid\ell^{3h}\mid\ell^{yh}$;
finally $\#J+\#s+1\le t\#s+2y^{20}\le(t+2)y^{20}\le y^{21}$, so $(\#J+\#s+1)\,yh\le y^{22}h$. If $h=0$ the bounds are trivial.
\end{proof}

\begin{theorem}[Triangle cost of Tate twists]\label{thm:triangle-cost}
Let the $w_j$ form a triangle of primes for $g_0$ up to $\ell^h$, $e<t$, and $Z$ a module of cocycles with values in $(O/\ell^\nu)(-e)$, unramified off $\ell$ and the $w_j$, $\bar w_j$, such that $\mathrm{cyc}(g_0)^{-e}-1$ kills the value at
$g_0$ of each member unramified off $\ell$. Let $b$ bound the length of $Z$ modulo the functions obeying the crystalline relations above $\ell$, plus the length of the odd classes. Then $Z$ has triangle cost $b$ along the $j$
of type $e$, up to $\ell^h$.
\end{theorem}

\begin{proof}
Put $\tau'_j=(\tau_j^c)^{-1}$ and $\varphi'_j=(\varphi_j^c)^{-1}$, and take for the set of the triangle cost the local set $\Lambda$ above $\ell$. \emph{Values at the inverse conjugates.} The $\tau_j$ and $\varphi_j$ act as the identity on $(O/\ell^\nu)(-e)$
($\mathrm{cyc}(\tau_j)=1$, $\mathrm{cyc}(\varphi_j)\equiv1$ modulo $\ell^\nu$), and where the action is trivial a cocycle has $z(g^{-1})=-z(g)$; so $z(\tau'_j)=-z(\tau_j^c)$ and $z(\tau_j)=z(\tau'_j)$ exactly when the even part $z+z\circ(g\mapsto g^c)$ vanishes at
$\tau_j$, and likewise for $\varphi_j$.

(a) Let $j$ have type $e$ (so $e>0$) and let $z\in Z$ vanish on $\Lambda$ and agree at the earlier pairs. By \cref{lem:even-parts}(a) its even part is an even class of exponent $e$ for the places $w_j$; it vanishes at the earlier $\tau_i$ and
at the earlier $\varphi_i$ of exponent $e$. By the triangle, $\ell^h$ kills its value at $\tau_j$, so $\ell^hz(\tau_j)=\ell^hz(\tau'_j)$.

(b) \emph{The first term.} A cocycle obeying the crystalline relations above $\ell$ vanishes on $\Lambda$: conjugating $g\in\Lambda$ by some $s$ into the decomposition group of the chosen place above its prime, crystalline cocycles vanish at
$sgs^{-1}$, and computing $z(s^{-1}\cdot sgs^{-1})=z(gs^{-1})$ in two ways gives $z(g)=0$. Hence the length of $Z$ modulo the functions vanishing on $\Lambda$ is at most its length modulo those obeying the crystalline relations.
\emph{The second term.} Let $S'\le Z$ consist of members vanishing on $\Lambda$, at all $\tau_j$, $\tau'_j$, and agreeing at $\varphi_j$, $\varphi'_j$ for $\mathrm{ty}(j)=e$. Its members are unramified off $\ell$ (\cref{lem:even-parts}(b)). If $e=0$ there is no sign condition and $S'$
lies in the odd classes. If $e\ne0$, the map $z\mapsto z\circ(g\mapsto g^c)-z$ sends $S'$ into the odd classes. Its kernel consists of even $z$; for such $z$ the agreement at $\varphi_j$, $\varphi'_j$ gives $2z(\varphi_j)=0$, and for each test point $g$
some $j$ of exponent $e$ has $z(\varphi_j)=2z(g)$, so $4z(g)=0$. With $d=\mathrm{cyc}(g_0)^{-e}-1$, which kills $z(g_0)$, the cocycle $4dz$ vanishes at $g_0$ and at the test points, so it is zero; and $4d$ divides $4d^2$, which divides
$\ell^h$. Hence the kernel is killed by $\ell^h$, and the length of $S'$ modulo its $\ell^h$-torsion is at most the length of the odd classes. The two terms add up to at most $b$.
\end{proof}
